% Generated draft: requires mathematical and source review.
\documentclass[11pt]{article}
\usepackage[margin=25mm]{geometry}
\usepackage{fontspec}
\setmainfont{Latin Modern Roman}
\usepackage{amsmath,amssymb,mathtools,unicode-math}
\setmathfont{Latin Modern Math}
\usepackage{enumitem,needspace,xurl,hyperref}
\setlength{\parindent}{0pt}
\setlength{\parskip}{5pt}
\setlength{\emergencystretch}{3em}
\newcommand{\N}{\mathbb{N}}
\newcommand{\Z}{\mathbb{Z}}
\newcommand{\Q}{\mathbb{Q}}
\newcommand{\R}{\mathbb{R}}
\newcommand{\C}{\mathbb{C}}
\newcommand{\F}{\mathbb{F}}
\newcommand{\A}{\mathbb{A}}
\newcommand{\PP}{\mathbb{P}}
\newcommand{\E}{\mathbb{E}}
\newcommand{\eps}{\varepsilon}
\newcommand{\dd}{\,\mathrm{d}}
\begin{document}
\section*{UnsolvedMath ID: 30000503}
\textbf{Source problem code:} catalogue entry
\subsection{Definitions and mathematical statement}
Determine the source of the difficulty of the regular inverse Galois problem: is it fundamentally diophantine rather than group-theoretic?

For background, a regular realization of a finite group $G$ over a field $k$ is a finite Galois extension $E/k(t)$ with
\[
 \operatorname{Gal}(E/k(t))\cong G,
 \qquad E\cap\overline{k}=k
\]
in a common algebraic closure, so that $k$ is the full algebraic constant field of $E$. The regular inverse Galois problem over $k$ asks whether every finite group has such a realization. The supplied record does not specify $k$; in particular, a rational-number-field formulation cannot be silently assumed from the name alone.

The actual source question asks for an explanation or mathematical analysis of the obstacles to this problem, not merely its yes-or-no resolution. It suggests that arithmetic questions about rational points or related diophantine conditions may be more fundamental than group-theoretic construction issues. The record supplies no definition of ``fundamentally,'' no class of allowed reductions, no separation of the proposed obstacle types, and no particular obstruction to prove necessary or sufficient. Consequently it does not determine a unique formal theorem comparing diophantine and group-theoretic difficulty.

The assessment mentions modular towers and Hurwitz spaces as methods producing special-case results. These provide context for the conceptual question but do not specify a universal moduli construction, rational-point criterion, or equivalence that can be stated here. No such criterion is invented, and no proof of the full regular inverse Galois problem is substituted for the source's requested diagnosis.
\subsection{Short English statement}
The entry asks whether the main obstacles to regular inverse Galois realization are diophantine in nature. It is a conceptual research question without a specified ground field or a formal test separating arithmetic and group-theoretic difficulty.
\subsection{Sources}
\begin{itemize}
\item Local UnsolvedMath registry, permanent ID 30000503, code catalogue entry, ``Diophantine Obstructions in the Regular Inverse Galois Problem.'' Proposed year 2006; proposer unspecified.
\item \emph{Pro-$p$ Extensions of Global Fields and pro-$p$ Groups} (2007), Oberwolfach Reports, pp.\ 1463--1536, DOI 10.4171/owr/2006/25; \url{https://doi.org/10.4171/owr/2006/25}. The displayed report year 2007 differs from the DOI and proposed year 2006; escaped dollar signs in the title are normalized as mathematical $p$.
\item Supplied ``Fried on modular towers and the regular inverse Galois problem'': \url{https://arxiv.org/abs/2008.01603}.
\item Supplied ``current RIGP overview'': \url{https://www.math.uci.edu/~mfried/deflist-cov/RIGP.html}.
\item Statement status exact; the supplied verification says the wording was confirmed without textual change.
\item Literature triage checked 2026-08-21: open. It describes the general regular inverse Galois problem as open and modular-tower/Hurwitz-space results as substantial special cases that do not answer the conceptual general question. The ground-field scope and current claims require independent verification. Created 2026-08-25; updated 2026-09-16.
\item Registry curation by Ulam AI / UnsolvedMath Contributors, CC BY 4.0; underlying report and literature provenance remain as supplied. Unreviewed statement-only draft.
\end{itemize}
\end{document}
